\section{Localization and Statistical Separation}
\label{arch:sec:localization}

\begin{proposition}[Localization of new violations]
\label{arch:prop:localization}
If $P_0$ is fair on $G$ and $\Impact$ is distributionally sound, every fairness
violation newly introduced by $P_1$ lies in $\ImpEdges$.
\end{proposition}

\begin{proof}
For $e=(u,v)\notin\ImpEdges$, both endpoints lie outside $\Impact$.  Soundness
implies that both endpoint output distributions, and therefore both expectations,
are unchanged.  Hence $g_{1,e}=g_{0,e}$, and the previous fairness certificate for
$P_0$ remains valid on $e$.
\end{proof}

A stochastic tester cannot in general distinguish an exact threshold from an
arbitrarily small violation with a uniform finite query bound.

\begin{proposition}[No finite uniform sharp-threshold test]
\label{arch:prop:zero-margin}
Fix $\tau\in(0,1/2)$ and $\delta\in(0,1/2)$.  On a graph containing one edge with
Bernoulli outputs, there is no tester with finite worst-case query bound that
uniformly distinguishes
\[
|p_u-p_v|\le\tau
\quad\text{from}\quad
|p_u-p_v|>\tau
\]
with success probability at least $1-\delta$ on every instance.
\end{proposition}

\begin{proof}
Take a fair instance on the boundary, e.g. $p_u=1/4$ and
$p_v=1/4+\tau$, and compare it with an unfair instance obtained by replacing
$p_v$ by $p_v+\eta$.  For any fixed finite query budget, the KL divergence between
the two possible transcript distributions tends to zero as $\eta\to0$.  Pinsker's
inequality therefore makes their total-variation distance arbitrarily small, so
no decision rule can maintain error at most $\delta$ in both worlds.
\end{proof}

Accordingly, the core problem uses a positive separation margin $\gamma>0$.
